\section{YAML}
\label{app:yaml}

\YAML{} is a human-friendly, structured data format designed to represent
\emph{configuration}, \emph{content}, and \emph{hierarchical data} in a
clear and readable way. %
Unlike \JSON{}, \YAML{} emphasises simplicity: its the indentation that
conveys the structure, and the syntax avoids unnecessary punctuation. %
Because of its readability and flexibility, \YAML{} is commonly used as a
data source for templating systems, configuration files, and structured
documents. %

\medskip

A \YAML{} file is composed of key--value pairs, lists, and nested
objects. %
Indentation determines the hierarchy, and the syntax supports a wide range
of scalar types (strings, numbers, booleans), collections, and optional
type annotations. %
Some common constructs include: %

\begin{itemize}
  \item \textbf{Mappings} (objects, or maps, or dictionnaries), e.g.:
        \yml{title: "My Document"} ;
  \item \textbf{Sequences} (lists), introduced with a dash, e.g.:
    \begin{ymlblock}
      - first item
      - second item
      - third item
    \end{ymlblock}
  \item \textbf{Nested structures}, created by indentation, e.g.:
    \begin{ymlblock}
      author:
          name: Alice
          id: 42
    \end{ymlblock}
\end{itemize}

\medskip

Of course, these structures can be combined and nested. %

\YAML’s design makes it particularly suitable as an external data source
for templating engines such as \seesection{app:liquid-templates}{Liquid}. %
Each \YAML{} key becomes accessible as a variable in the template, and the
structure of the \YAML{} document (lists, maps, nested objects) is
preserved during rendering. %


%%% Local Variables:
%%% mode: LaTeX
%%% TeX-master: "../liquidtex"
%%% End:
